\section{Boundary algebra from screens (holographic boundary object)}
\label{app:holo_boundary_algebra_from_screens}

\paragraph{Purpose.}
\AuditTag This appendix instantiates the boundary object of Appendix~\ref{app:holo_scope_contract}:
it fixes a concrete region family on the protocol screen and defines the corresponding boundary algebras
as a restriction of the protocol-induced local net.

\subsection{Holo1: screen regions as a concrete region family}
\label{subsec:holo1_screen_regions}

\paragraph{Screen and regions.}
\InterfaceTag When spatial statistics are invoked, the manuscript fixes an addressing basis and a finite screen
at order $n$ (Section~\ref{sec:hilbert_addressing}).
Let $S_n$ denote the finite site set on the chosen addressing screen at order $n$
(e.g.\ a $2^n\times 2^n$ grid for the 2D Hilbert screen).
A \emph{screen region} is a subset $R\subseteq S_n$.

\paragraph{Region-to-window association.}
\AuditTag A screen region $R$ should be viewed as a quotient class of concrete finite windows whose readouts
are supported on $R$.
This matches the abstract region construction $\mathfrak{R}_{\mathrm{prot}}=\mathcal{W}/\sim$
in Appendix~\ref{app:construct_local_net_from_protocol}.

\subsection{Holo1: definition of the boundary algebra}
\label{subsec:holo1_boundary_algebra_def}

\begin{definition}[Boundary algebra on a screen region]
\label{def:holo_boundary_algebra}
\MathTag Fix a protocol state $(m,n,K)$ and a screen site set $S_n$.
For a screen region $R\subseteq S_n$, define the \emph{boundary algebra}
\[
\mathcal{A}_{\partial}(R)\ :=\ \mathcal{A}_{\mathrm{prot}}(\mathcal{O}_R),
\]
where $\mathcal{O}_R\in\mathfrak{R}_{\mathrm{prot}}$ is the protocol region class represented by windows supported on $R$
in the operational inclusion sense (Appendix~\ref{app:construct_local_net_from_protocol}).
\end{definition}

\paragraph{Concrete finite models.}
\InterfaceTag At fixed finite resolution, one may realize $\mathcal{A}_{\partial}(R)$ in either of the two
canonical finite-resolution models of Appendix~\ref{app:observable_algebra_state_base}:
\begin{itemize}
\item \textbf{Classical readout:} $\mathcal{A}_{\partial}(R)=\mathcal{A}_{\mathrm{cl}}(\Omega_R)=\{f:\Omega_R\to\CC\}$,
for a finite outcome set $\Omega_R$.
\item \textbf{Quantum readout:} $\mathcal{A}_{\partial}(R)\cong \mathcal{B}(\mathcal{H}_R)\cong M_{d(R)}(\CC)$
for a finite-dimensional Hilbert space $\mathcal{H}_R$ representing a declared local register on $R$.
\end{itemize}
The manuscript treats the chosen realization as part of the declared interface/modeling layer (ledger-scoped).

\subsection{Holo1: PT properties inherited from the protocol-induced net}
\label{subsec:holo1_pt_properties}

\begin{proposition}[Isotony of the boundary algebra]
\label{prop:holo_boundary_isotony}
\MathTag If $R_1\subseteq R_2\subseteq S_n$ (as subsets of sites), then
\[
\mathcal{A}_{\partial}(R_1)\subseteq \mathcal{A}_{\partial}(R_2)
\]
as a unital $^\ast$-subalgebra inclusion under the inductive-limit embeddings.
\end{proposition}
\begin{proof}
This is the restriction of Proposition~\ref{prop:prot_net_isotony} to the screen-induced region family:
operational inclusion $R_1\subseteq R_2$ induces $\mathcal{O}_{R_1}\subseteq \mathcal{O}_{R_2}$ in $\mathfrak{R}_{\mathrm{prot}}$.
\end{proof}

\begin{corollary}[Microcausality in the tensor-product subclass (PT carrier)]
\label{cor:holo_boundary_microcausality_tensor_subclass}
\MathTag In the tensor-product readout subclass (Definition~\ref{def:tensor_product_readout_subclass}),
the boundary net $R\mapsto\mathcal{A}_{\partial}(R)$ satisfies microcausality for operationally disjoint regions.
\end{corollary}
\begin{proof}
This is immediate from Corollary~\ref{cor:microcausality_tensor_subclass} applied to the restricted region family.
\end{proof}

\subsection{Holo1: identifiability boundary (what boundary summaries can and cannot encode)}
\label{subsec:holo1_identifiability_boundary}

\paragraph{Class-function summaries are informationally bounded.}
\AuditTag Many protocol diagnostics report only class-function loop summaries (e.g.\ cycle-type histograms for $S_4$ holonomy).
In such cases, the boundary data is subject to the identifiability boundary:
two realizations with identical cycle-type histograms are indistinguishable by all class-function observables
(Proposition~\ref{prop:cycle_type_identifiability_boundary}).
Any holographic claim that requires distinguishing such realizations must therefore enlarge the observable family
or be downgraded (Failure point H0c of Appendix~\ref{app:holo_scope_contract}).

\paragraph{Deterministic sample fragment (region chain).}
\AuditTag The following generated fragment instantiates a small nested region chain and its protocol-capacity summaries,
as an auditable companion to isotony:
\begin{table}[H]
\centering
\scriptsize
\setlength{\tabcolsep}{6pt}
\renewcommand{\arraystretch}{1.15}
\begin{tabular}{c c c c}
\toprule
region & $|R|$ & $I_{\mathrm{prot}}(m,n;R)$ (bits) & included in \\
\midrule
R_1 & 1 & 6 & - \\
R_2 & 4 & 24 & R_1 \\
R_3 & 16 & 96 & R_2 \\
R_4 & 64 & 384 & R_3 \\
\bottomrule%
\end{tabular}
\caption{Boundary-region sample (nested chain) for isotony bookkeeping. Rows are generated by \texttt{scripts/exp\_holo\_boundary\_algebra\_sample.py}.}
\label{tab:holo_boundary_algebra_sample}
\end{table}
\paragraph{Boundary-region sample (audit).} \AuditTag This fragment instantiates a nested region chain by site counts (R$_1\subset$R$_2\subset$R$_3\subset$R$_4$) at anchor (m,n)=(6,3). The protocol capacity $I_{\mathrm{prot}}(m,n;R)=m|R|$ is monotone along the chain, consistent with isotony of the associated boundary algebras (Appendix~\ref{app:holo_boundary_algebra_from_screens}).%
